% 自动生成：python paper/build_tables.py（勿手工编辑）
\begin{table}[htbp]
\centering
\caption{预热阶段不同时刻的温度分布（单位：℃）}
\label{tab:t1}
\begin{tabular}{lccccc}
\toprule
时间/s & 0 cm & 0.5 cm & 1 cm & 1.5 cm & 2 cm \\
\midrule
100 & 28.0001 & 28.0003 & 28.0040 & 28.0327 & 28.1801 \\
300 & 28.0408 & 28.0635 & 28.1514 & 28.3681 & 28.8489 \\
600 & 28.4533 & 28.5360 & 28.8039 & 29.3159 & 30.1652 \\
900 & 29.3243 & 29.4583 & 29.8755 & 30.6161 & 31.7304 \\
1200 & 30.5427 & 30.7098 & 31.2223 & 32.1126 & 33.4276 \\
1500 & 31.9957 & 32.1867 & 32.7660 & 33.7463 & 35.1203 \\
1800 & 33.5753 & 33.7720 & 34.3642 & 35.3621 & 36.7856 \\
\bottomrule
\end{tabular}
\end{table}

\begin{table}[htbp]
\centering
\caption{预热阶段不同时刻的干基含水率分布（单位：$\mathrm{kg\,kg^{-1}}$）}
\label{tab:t2}
\begin{tabular}{lccccc}
\toprule
时间/s & 0 cm & 0.5 cm & 1 cm & 1.5 cm & 2 cm \\
\midrule
100 & 2.5500 & 2.5500 & 2.5500 & 2.5500 & 2.2470 \\
300 & 2.5500 & 2.5500 & 2.5500 & 2.5492 & 2.0508 \\
600 & 2.5500 & 2.5500 & 2.5500 & 2.5353 & 1.8770 \\
900 & 2.5500 & 2.5500 & 2.5497 & 2.5045 & 1.7547 \\
1200 & 2.5500 & 2.5500 & 2.5482 & 2.4646 & 1.6586 \\
1500 & 2.5500 & 2.5499 & 2.5445 & 2.4206 & 1.5787 \\
1800 & 2.5500 & 2.5497 & 2.5383 & 2.3755 & 1.5102 \\
\bottomrule
\end{tabular}
\end{table}

\begin{table}[htbp]
\centering
\caption{烘干前 3 小时不同时刻的温度分布（单位：℃）}
\label{tab:t3}
\begin{tabular}{lccccc}
\toprule
时间/h & 0 cm & 0.5 cm & 1 cm & 1.5 cm & 2 cm \\
\midrule
0 & 28.0000 & 28.0000 & 28.0000 & 28.0000 & 28.0000 \\
0.5 & 32.1893 & 32.3821 & 32.9659 & 33.9608 & 35.4130 \\
1 & 40.3816 & 40.5536 & 41.0605 & 41.8782 & 42.9976 \\
1.5 & 45.8468 & 45.9348 & 46.1929 & 46.6051 & 47.1400 \\
2 & 48.4502 & 48.4881 & 48.5984 & 48.7735 & 49.0033 \\
2.5 & 49.4670 & 49.4792 & 49.5136 & 49.5653 & 49.6609 \\
3 & 49.8495 & 49.8553 & 49.8746 & 49.9101 & 49.9664 \\
\bottomrule
\end{tabular}
\end{table}

\begin{table}[htbp]
\centering
\caption{烘干前 3 小时不同时刻的干基含水率分布（单位：$\mathrm{kg\,kg^{-1}}$）}
\label{tab:t4}
\begin{tabular}{lccccc}
\toprule
时间/h & 0 cm & 0.5 cm & 1 cm & 1.5 cm & 2 cm \\
\midrule
0 & 2.5500 & 2.5500 & 2.5500 & 2.5500 & 2.5500 \\
0.5 & 2.5499 & 2.5489 & 2.5256 & 2.3257 & 1.6486 \\
1 & 2.5257 & 2.4948 & 2.3578 & 2.0230 & 1.4711 \\
1.5 & 2.3861 & 2.3256 & 2.1344 & 1.8020 & 1.3476 \\
2 & 2.1708 & 2.1084 & 1.9236 & 1.6259 & 1.2311 \\
2.5 & 1.9566 & 1.9006 & 1.7360 & 1.4720 & 1.1166 \\
3 & 1.7662 & 1.7165 & 1.5702 & 1.3333 & 1.0081 \\
\bottomrule
\end{tabular}
\end{table}

\begin{table}[htbp]
\centering
\caption{固定半径工况代表性时刻的干基含水率（单位：$\mathrm{kg\,kg^{-1}}$）}
\label{tab:t5}
\begin{tabular}{lccccc}
\toprule
时间/h & 0 cm & 0.5 cm & 1 cm & 1.5 cm & 2 cm \\
\midrule
0 & 2.5500 & 2.5500 & 2.5500 & 2.5500 & 2.5500 \\
6 & 1.0156 & 0.9868 & 0.9006 & 0.7546 & 0.5335 \\
12 & 0.4548 & 0.4418 & 0.4019 & 0.3289 & 0.1640 \\
18 & 0.2979 & 0.2906 & 0.2679 & 0.2247 & 0.0842 \\
24 & 0.2371 & 0.2321 & 0.2161 & 0.1851 & 0.0663 \\
30 & 0.2053 & 0.2013 & 0.1887 & 0.1637 & 0.0597 \\
36 & 0.1854 & 0.1821 & 0.1714 & 0.1501 & 0.0565 \\
42 & 0.1717 & 0.1688 & 0.1594 & 0.1405 & 0.0547 \\
48 & 0.1615 & 0.1589 & 0.1504 & 0.1332 & 0.0536 \\
54 & 0.1536 & 0.1511 & 0.1434 & 0.1275 & 0.0528 \\
烘干结束时间（57.1904 h） & 0.1500 & 0.1477 & 0.1402 & 0.1249 & 0.0525 \\
\bottomrule
\end{tabular}
\parbox{0.88\linewidth}{\footnotesize\textit{注：}表中每 6 h 取样；末行为连续事件定位得到的精确烘干终止时刻，不是 60 s 规则输出时间点。}
\end{table}

\begin{table}[htbp]
\centering
\caption{收缩工况代表性时刻的干基含水率（单位：$\mathrm{kg\,kg^{-1}}$）}
\label{tab:t6}
\begin{tabular}{lccccc}
\toprule
时间/h & 0 cm & 0.5 cm & 1 cm & 1.5 cm & 药材表面 \\
\midrule
0 & 2.5500 & 2.5500 & 2.5500 & 2.5500 & 2.5500 \\
6 & 1.7164 & 1.5350 & 1.0212 &  & 0.4215 \\
12 & 0.7340 & 0.6516 & 0.4068 &  & 0.1666 \\
18 & 0.4065 & 0.3667 & 0.2387 &  & 0.0888 \\
24 & 0.2837 & 0.2598 & 0.1783 &  & 0.0671 \\
30 & 0.2253 & 0.2085 & 0.1488 &  & 0.0593 \\
36 & 0.1920 & 0.1790 & 0.1313 &  & 0.0558 \\
42 & 0.1706 & 0.1598 & 0.1196 &  & 0.0539 \\
48 & 0.1556 & 0.1463 & 0.1113 &  & 0.0529 \\
烘干结束时间（50.8235 h） & 0.1500 & 0.1413 & 0.1081 &  & 0.0525 \\
\bottomrule
\end{tabular}
\parbox{0.88\linewidth}{\footnotesize\textit{注：}表中每 6 h 取样；0--1.5 cm 为固定物理坐标，空白表示该位置已位于收缩后药材外部；“药材表面”为移动坐标 $r=R(t)$。末行为连续事件定位得到的精确烘干终止时刻。}
\end{table}
